\documentclass[11pt,a4paper]{article}
\usepackage[utf8]{inputenc}
\usepackage[T1]{fontenc}
\usepackage[ngerman]{babel}
\usepackage[margin=2.5cm]{geometry}
\usepackage{amsmath,amssymb}
\usepackage{listings}
\usepackage{xcolor}
\usepackage{parskip}

\definecolor{codegray}{gray}{0.95}
\definecolor{kaminrot}{HTML}{9B111E}
\usepackage{sectsty}
\allsectionsfont{\color{kaminrot}}

\lstset{
  language=Python,
  basicstyle=\ttfamily\small,
  backgroundcolor=\color{codegray},
  keywordstyle=\color{kaminrot}\bfseries,
  commentstyle=\color{gray}\itshape,
  breaklines=true,
  frame=single,
  framerule=0pt,
  xleftmargin=0.5em,
  aboveskip=1em,
  belowskip=1em
}

\title{\color{kaminrot}Per-Pulver-Auswertung von BH-Loops\\
\large Losses und $B_\mathrm{peak}$ pro Pulversorte in magnum.np}
\author{}
\date{17.\,Juli 2026}

\begin{document}
\maketitle

\noindent\textbf{Setup:} Pulver-Simulation mit gro"sen Partikeln (4\,\textmu m) auf fcc-Gitterpl"atzen und kleinen Partikeln (1\,\textmu m) in den Zwickeln, Mischungsverh"altnis 70:30. BH-Loop mit AC-Feld in $x$-Richtung. Gesucht: Losses und $B_\mathrm{peak}$ pro Pulversorte.

\section{Was \texttt{average(m $\cdot$ Ms $\cdot$ filter)} wirklich liefert}

Der \texttt{ScalarLogger} schickt jeden Tensor durch \texttt{state.avg()}, und das ist auf "aquidistanten Meshes ein Mittelwert "uber \emph{alle} Zellen des Meshes (\texttt{state.py:159}, \texttt{scalar\_logger.py:105}). Man erh"alt also nicht die mittlere Magnetisierung \emph{im} Pulver $i$, sondern deren Beitrag zum Gesamtvolumen:
\begin{equation}
\langle M_s\, m\, f_i \rangle_\text{mesh} \;=\; \varphi_i \, M_i
\qquad\text{mit } \varphi_i = \text{Volumenanteil der Phase } i .
\end{equation}
F"ur per-Pulver-Gr"o"sen muss durch \texttt{state.avg(filter\_i)} geteilt werden.

\textbf{Sanity-Checks:}
\begin{enumerate}
  \item Das phasennormierte $J_i^\text{peak} = \mu_0 M_i^\text{peak}$ muss $\le \mu_0 M_s$ der jeweiligen Phase sein. Wenn nicht, stimmt Filter oder Normierung nicht.
  \item Ohne Normierung w"are das kleine Pulver durch den 30\,\%-Anteil ohnehin um Faktor ${\sim}2{,}3$ unterdr"uckt. Wenn es trotzdem gr"o"ser aussieht: zuerst die Filter pr"ufen.
\end{enumerate}

\section{Ben"otigte Outputs}

Zus"atzlich zum bisherigen Output braucht man das \textbf{gefilterte, phasennormierte Entmagnetisierungsfeld}:

\begin{lstlisting}
mu0 = constants.mu_0
phi_L = state.avg(filter_L)   # Volumenanteile, einmal berechnen
phi_S = state.avg(filter_S)

logger = ScalarLogger("data/loop.dat", [
    't', external.h,                                # H_ext(t)
    ('J_L', lambda s: mu0 * s.material["Ms"] * s.m * filter_L / phi_L),
    ('J_S', lambda s: mu0 * s.material["Ms"] * s.m * filter_S / phi_S),
    ('h_L', lambda s: (demag.h(s) + external.h(s)) * filter_L / phi_L),
    ('h_S', lambda s: (demag.h(s) + external.h(s)) * filter_S / phi_S),
])
\end{lstlisting}

Daraus in der Auswertung:

\paragraph{$B$ pro Pulver.} Mit dem \emph{lokalen} Feld:
\begin{equation}
B_i \;=\; \mu_0 \langle H \rangle_i + J_i ,
\qquad
\langle H \rangle_i = H_\text{ext} + \langle H_\text{demag} \rangle_i .
\end{equation}
$\mu_0 (H_\text{ext} + M_i)$ ist \emph{nicht} die Flussdichte im Partikel: Im weichmagnetischen Partikel ist das Entmagnetisierungsfeld gro"s (Gr"o"senordnung $-M/3$ f"ur die Kugel, plus Streufelder der Nachbarn) und f"ur beide Sorten verschieden. Das kann die Reihenfolge der $B_\mathrm{peak}$-Werte deutlich verschieben.

\paragraph{Losses pro Pulver.} Ringintegral "uber genau eine Periode (nach dem Einschwingen), numerisch per Trapezregel:
\begin{equation}
w_i \;=\; \oint H_x \, \mathrm{d}J_{i,x}
\qquad\text{(pro Einheit Phasenvolumen, pro Zyklus)},
\end{equation}
absolute Verluste: $W_i = w_i \cdot \varphi_i V$.

\section{$H_\text{ext}$ oder Streufelder dazu?}

Beides ist \glqq richtig\grqq, aber f"ur verschiedene Fragen:

\paragraph{(i) Gesamt-Loop / Gesamtverluste.} $H_\text{ext}$ ist \emph{exakt} richtig:
\begin{equation}
w_\text{tot} \;=\; \mu_0 \oint H_\text{ext} \, \mathrm{d}M_\text{tot} .
\end{equation}
"Uber einen geschlossenen Zyklus ist die Streufeldenergie eine Zustandsfunktion und tr"agt netto nichts bei.

\paragraph{(ii) Per-Pulver mit $H_\text{ext}$.}
\begin{equation}
w_i \;=\; \mu_0 \oint H_\text{ext} \, \mathrm{d}M_i
\end{equation}
ist die \emph{vom externen Feld in Phase $i$ eingebrachte Arbeit}. Die Summe ergibt exakt den Gesamtverlust (guter Konsistenz-Check!):
\begin{equation}
\varphi_L\, w_L + \varphi_S\, w_S \;=\; w_\text{tot} .
\end{equation}
Aber: Diese Zerlegung ignoriert magnetostatischen Energietransfer zwischen den Sorten. Energie, die "uber das Feld des gro"sen Pulvers ins kleine gepumpt (oder umgekehrt) und \emph{dort} dissipiert wird, landet in der falschen Spalte.

\paragraph{(iii) Per-Pulver tats"achliche Dissipation.} Daf"ur braucht man das lokale Feld:
\begin{equation}
w_i^\text{dis} \;\approx\; \oint \langle H \rangle_{i,x} \, \mathrm{d}J_{i,x} ,
\qquad
\langle H \rangle_i = H_\text{ext} + \langle H_\text{demag} \rangle_i .
\end{equation}
Dabei einfach das volle \texttt{demag.h} filtern, nicht nach Quelle trennen: Der Eigenanteil hebt sich "uber den geschlossenen Zyklus n"aherungsweise weg, "ubrig bleibt effektiv der Streufeldbeitrag der jeweils anderen Sorte. Kleine Einschr"ankung: $\langle H \rangle_i \, \mathrm{d}\langle M \rangle_i$ ist das Produkt der Mittelwerte statt des Mittelwerts des Produkts, also eine N"aherung --- f"ur die Frage \glqq welche Sorte dissipiert wie viel\grqq{} aber der pragmatische Standardweg.

Die Differenz zwischen (ii) und (iii) ist gerade der Energietransfer zwischen den Pulvern --- die auszurechnen lohnt sich, bevor man das Ranking interpretiert.

\paragraph{Muss der Eigenanteil aus dem Demagfeld herausgerechnet werden?}
Nein --- und die Begr"undung ist \emph{nicht}, dass er klein w"are (das ist er nicht: im Partikel von der Gr"o"senordnung $-M/3$). Entscheidend ist, dass sein Beitrag zum Ringintegral "uber einen geschlossenen Zyklus exakt verschwindet, weil die Selbstenergie eine Zustandsfunktion der eigenen Konfiguration ist. Mit
\begin{equation}
E_i^\text{self} \;=\; -\frac{\mu_0}{2} \int_{V_i} H_\mathrm{d}[M_i] \cdot M_i \,\mathrm{d}V
\end{equation}
gilt wegen der Symmetrie des Demag-Operators
\begin{equation}
\mu_0 \oint \int_{V_i} H_\mathrm{d}[M_i] \cdot \mathrm{d}M \,\mathrm{d}V
\;=\; -\oint \mathrm{d}E_i^\text{self} \;=\; 0
\qquad \text{(exakt, da die Konfiguration zur"uckkehrt).}
\end{equation}
Der Eigenanteil ist also konservativ, nicht klein --- er leistet "uber den Zyklus netto keine Arbeit.

\emph{Warum leistet das Self-Feld netto keine Arbeit, obwohl $M$ Hysterese zeigt?} Hysterese ist eine Doppeldeutigkeit von $M$ gegen"uber $H_\text{ext}$: Bei demselben $H_\text{ext}$ sind auf Auf- und Abw"artsast verschiedene $M$ m"oglich, weil $H_\text{ext}$ eine unabh"angige, von au"sen aufgepr"agte Gr"o"se ist. Das Self-Feld ist dagegen eine \emph{eindeutige} Funktion der momentanen Konfiguration --- dieselbe Magnetisierungsverteilung erzeugt immer dasselbe Self-Feld. Es wirkt wie eine Feder an der Magnetisierung: Auf dem aufsteigenden Ast wird Arbeit gegen das Self-Feld geleistet und als Streufeldenergie \emph{gespeichert}, auf dem absteigenden Ast kommt genau diese Energie zur"uck. Formal ist $H_\mathrm{d}[M]\cdot\mathrm{d}M$ ein totales Differential der Selbstenergie (Gl.~9), $H_\text{ext}\cdot\mathrm{d}M$ dagegen nicht --- deshalb hat die eine Schleife zwingend Fl"ache null, die andere nicht. Die Verlustenergie stammt aus der D"ampfung: An Instabilit"aten (Nukleation, Depinning) springt die Magnetisierung dynamisch ins n"achste Minimum, und die Energiedifferenz wird "uber den $\alpha$-Term der LLG zu W"arme. Das Demag-Feld formt die Energielandschaft mit --- es bestimmt, \emph{wie viel} dissipiert wird --- tr"agt aber als Zustandsfunktion zum Ringintegral netto nichts bei. Voraussetzung ist, dass die \emph{gesamte} Konfiguration zur"uckkehrt (Zelle f"ur Zelle, nicht nur der Mittelwert) --- im eingeschwungenen AC-Zyklus erf"ullt; daher die letzte volle Periode auswerten, nicht die Neukurve.

Eine N"aherung entsteht erst durch die Phasen\emph{mittelung}: Das Integral
$\oint \langle H^\text{self} \rangle_i \, \mathrm{d}\langle M \rangle_i$
verschwindet nur dann exakt, wenn $\langle H^\text{self} \rangle_i$ eine eindeutige Funktion von $\langle M \rangle_i$ ist. Im Mean-Field-Bild,
\begin{equation}
\langle H^\text{self} \rangle_i \;\approx\; -N_\text{eff} \, \langle M \rangle_i ,
\end{equation}
ist die R"uckstellung linear und die eingeschlossene Fl"ache exakt null. Ein Restbeitrag bleibt, wenn die interne Dom"anenkonfiguration bei gleichem $\langle M \rangle$ auf Auf- und Abw"artsast verschieden ist --- dann ist die Kurve $\langle H^\text{self} \rangle$ vs.\ $\langle M \rangle$ selbst hysteretisch. \emph{Dieser Rest} ist die Kleinheitsannahme, typischerweise klein gegen die Loop-Fl"ache.

Wer es exakt will, hat zwei Optionen:
\begin{enumerate}
  \item \textbf{Quellentrennung "uber die Linearit"at des Demag-Operators:} $H_\mathrm{d}^\text{total} = H_\mathrm{d}[\text{nur gro"ses Pulver}] + H_\mathrm{d}[\text{nur kleines Pulver}]$. Das Kreuzfeld auf Phase $i$ erh"alt man, indem man das Demagfeld mit genulltem $M_s$ der eigenen Phase auswertet (eine zus"atzliche Demag-Auswertung pro Logpunkt, oder Postprocessing "uber gespeicherte $m$-Snapshots).
  \item \textbf{Lokale LLG-Dissipation direkt loggen:}
  \begin{equation}
  p_\text{dis} \;=\; \frac{\alpha \gamma \mu_0 M_s}{1+\alpha^2}\, \left| m \times H_\text{eff} \right|^2 ,
  \end{equation}
  gefiltert pro Phase. Das ist die exakte Dissipationsdichte der LLG-Dynamik und macht die Zuordnung eindeutig, ganz ohne Feld-Zerlegung.
\end{enumerate}

Wichtig: F"ur $B_i$ (Gl.~2) bleibt der Eigenanteil selbstverst"andlich enthalten --- die Wegheb-Eigenschaft betrifft nur das Verlust-Ringintegral, nicht die Flussdichte im Partikel.

\section{Exakte Quellentrennung im Logger}

Der Demag-Operator ist linear in $M_s\,m$ (\texttt{demagPBC.py:72} rechnet $\mathrm{FFT}(M_s\, m)$ mal Kernel), also gilt \emph{exakt}
\begin{equation}
H_\mathrm{d}^\text{total} \;=\; H_\mathrm{d}[\text{nur gro"ses Pulver}] \;+\; H_\mathrm{d}[\text{nur kleines Pulver}] .
\end{equation}
Es gen"ugt daher \emph{eine} zus"atzliche Demag-Auswertung mit maskierter Magnetisierung; der Rest ergibt sich als Differenz. \texttt{state.m} hat keinen speziellen Setter (kein Normalisieren, kein $M$-Caching im Feldterm --- nur der $N$-Tensor ist gecacht und h"angt nicht von $m$ ab), deshalb funktioniert tempor"ares Maskieren sauber:

\begin{lstlisting}
def demag_h_partial(state, filt):
    """Demagfeld, das nur von der Magnetisierung in filt erzeugt wird."""
    m_full = state.m
    try:
        state.m = m_full * filt      # nur Quellphase behalten
        return demag.h(state)
    finally:
        state.m = m_full             # Originalobjekt wiederherstellen

logger = ScalarLogger("data/loop.dat", [
    't', external.h,
    ('J_L',  lambda s: mu0 * s.material["Ms"] * s.m * filter_L / phi_L),
    ('J_S',  lambda s: mu0 * s.material["Ms"] * s.m * filter_S / phi_S),
    ('hd_L', lambda s: s.h_demag  * filter_L / phi_L),  # volles Feld in L
    ('hd_S', lambda s: s.h_demag  * filter_S / phi_S),
    ('hL_L', lambda s: s.h_from_L * filter_L / phi_L),  # Eigenfeld L
    ('hL_S', lambda s: s.h_from_L * filter_S / phi_S),  # Kreuzfeld L -> S
])

for i in tqdm(torch.arange(0, T, dt)):
    llg.step(state, dt)
    state.h_demag  = demag.h(state)                    # 1 FFT-Faltung
    state.h_from_L = demag_h_partial(state, filter_L)  # 1 weitere
    logger << state
\end{lstlisting}

Damit sind alle vier Anteile verf"ugbar, zwei davon per Subtraktion im Postprocessing (exakt bis auf Float-Rundung):
\begin{center}
\begin{tabular}{lll}
\textbf{Anteil} & \textbf{Quelle $\to$ Ort} & \textbf{Spalte} \\[2pt]
Eigenfeld im gro"sen Pulver   & L $\to$ L & \texttt{hL\_L} \\
Kreuzfeld im kleinen Pulver   & L $\to$ S & \texttt{hL\_S} \\
Kreuzfeld im gro"sen Pulver   & S $\to$ L & \texttt{hd\_L} $-$ \texttt{hL\_L} \\
Eigenfeld im kleinen Pulver   & S $\to$ S & \texttt{hd\_S} $-$ \texttt{hL\_S} \\
\end{tabular}
\end{center}

\textbf{Kosten:} eine zus"atzliche FFT-Faltung pro Logpunkt. Ein LLG-Schritt macht intern mehrere RHS-Auswertungen mit je einer Demag-Faltung, das f"allt also kaum ins Gewicht (mit \texttt{every=} im Logger vernachl"assigbar). \textbf{Alternative ohne Eingriff in den Run:} $m$-Snapshots als VTI speichern und die Partialfelder offline mit derselben Maskierungs-Idee berechnen --- f"ur glatte Ringintegrale braucht das aber ausreichende Zeitaufl"osung.

\textbf{Konsistenz-Check Energietransfer:} Mit den Kreuzfeldern l"asst sich der Energietransfer zwischen den Sorten direkt berechnen, $\oint \langle h_\text{cross}\rangle_i \,\mathrm{d}J_i$. Die beiden Transferterme m"ussen sich "uber den geschlossenen Zyklus (bis auf die Mittelungs-N"aherung aus Abschnitt~3) gegenseitig aufheben.

\section{Warum das Ergebnis trotzdem plausibel ist}

\glqq Losses h"angen stark von $B_\mathrm{peak}$ ab\grqq{} (Steinmetz-Denke) gilt f"ur \emph{ein} Material bei variierender Amplitude --- nicht beim Vergleich zweier verschiedener Mikrostrukturen. Die Loop-Fl"ache ist Amplitude $\times$ Breite, und die Breite unterscheidet sich hier systematisch:

\begin{itemize}
  \item Die 1-\textmu m-Partikel sitzen in den Zwickeln zwischen den gro"sen Kugeln und wirken als \textbf{Flussbr"ucken}: Der Fluss der gro"sen Nachbarn wird durch sie hindurchgeb"undelt $\Rightarrow$ lokal h"oheres $B$ (der beobachtete hohe $B_\mathrm{peak}$). Gleichzeitig sind sie durch die Streufelder der gro"sen Nachbarn stark vorgespannt: Ihre lokale Schleife ist geschert und fast reversibel $\Rightarrow$ kleine Fl"ache trotz gro"ser Amplitude.
  \item Die 4-\textmu m-Partikel sind gro"s genug f"ur Dom"anenw"ande und Vortex-Zust"ande. Sie schalten "uber Nukleation und Pinning $\Rightarrow$ breite Schleife $\Rightarrow$ h"ohere Verluste trotz kleinerem $B_\mathrm{peak}$.
\end{itemize}

Also kein Widerspruch --- aber in dieser Reihenfolge verifizieren:
\begin{enumerate}
  \item Normierung fixen (Abschnitt~1),
  \item $B$ mit lokalem Feld neu berechnen (Abschnitt~2),
  \item Losses einmal mit $H_\text{ext}$ und einmal mit $\langle H \rangle_i$ rechnen und pr"ufen, ob das Ranking in beiden Varianten bestehen bleibt.
\end{enumerate}

\end{document}
